% =============================================================================
% Appendix C: Debugging Guide
% =============================================================================
\chapter{Debugging Guide}
\label{app:debugging}

This appendix provides solutions to common issues encountered when training language models.

% -----------------------------------------------------------------------------
\section{NaN Loss}
% -----------------------------------------------------------------------------

\textbf{Symptoms}: Loss becomes \texttt{nan} or \texttt{inf} during training.

\subsection{Causes}

\begin{enumerate}
    \item \textbf{Learning rate too high}: Gradients explode
    \item \textbf{No gradient clipping}: Large gradients cause overflow
    \item \textbf{Numerical instability in attention}: $\exp$ overflow in softmax
    \item \textbf{Division by zero}: Usually in normalization layers
    \item \textbf{Bad data}: Infinite or NaN values in input
\end{enumerate}

\subsection{Solutions}

\begin{lstlisting}[style=python]
# 1. Reduce learning rate
learning_rate = 1e-4  # Instead of 3e-4

# 2. Enable gradient clipping
torch.nn.utils.clip_grad_norm_(model.parameters(), max_norm=1.0)

# 3. Check for NaN in gradients
for name, param in model.named_parameters():
    if param.grad is not None and torch.isnan(param.grad).any():
        print(f"NaN gradient in {name}")

# 4. Check data for NaN
assert not torch.isnan(input_ids).any(), "NaN in input"

# 5. Add epsilon to divisions
rms = torch.sqrt(torch.mean(x ** 2, dim=-1, keepdim=True) + 1e-6)
\end{lstlisting}

\subsection{Debugging Script}

\begin{lstlisting}[style=python]
def check_model_health(model, batch):
    """Check for numerical issues."""
    model.eval()
    with torch.no_grad():
        output = model(**batch)

        # Check loss
        if torch.isnan(output["loss"]) or torch.isinf(output["loss"]):
            print("WARNING: Loss is NaN or Inf")

        # Check logits
        logits = output["logits"]
        print(f"Logits - min: {logits.min():.2f}, max: {logits.max():.2f}")

        if torch.isnan(logits).any():
            print("WARNING: NaN in logits")

        # Check for extreme values
        if logits.abs().max() > 100:
            print("WARNING: Extreme logit values")
\end{lstlisting}

% -----------------------------------------------------------------------------
\section{Loss Not Decreasing}
% -----------------------------------------------------------------------------

\textbf{Symptoms}: Loss stays flat or decreases very slowly.

\subsection{Causes and Solutions}

\begin{enumerate}
    \item \textbf{Learning rate too low}:
    \begin{lstlisting}[style=python]
# Try increasing learning rate
learning_rate = 1e-3  # Instead of 1e-4
    \end{lstlisting}

    \item \textbf{Bug in data pipeline}:
    \begin{lstlisting}[style=python]
# Verify data looks correct
batch = next(iter(train_loader))
print(tokenizer.decode(batch["input_ids"][0]))
print(f"Labels: {batch['labels'][0][:20]}")
    \end{lstlisting}

    \item \textbf{Labels not shifted correctly}:
    \begin{lstlisting}[style=python]
# Correct shift operation
input_ids = tokens[:-1]
labels = tokens[1:]
    \end{lstlisting}

    \item \textbf{Model too small}:
    \begin{lstlisting}[style=python]
# Increase capacity
config = ModelConfig(
    n_layers=8,    # Was 6
    d_model=512,   # Was 384
)
    \end{lstlisting}

    \item \textbf{Optimizer not updating}:
    \begin{lstlisting}[style=python]
# Ensure backward and step are called
loss.backward()
optimizer.step()
optimizer.zero_grad()  # Don't forget this!
    \end{lstlisting}
\end{enumerate}

% -----------------------------------------------------------------------------
\section{Out of Memory (OOM)}
% -----------------------------------------------------------------------------

\textbf{Symptoms}: \texttt{CUDA out of memory} error.

\subsection{Solutions}

\begin{enumerate}
    \item \textbf{Reduce batch size}:
    \begin{lstlisting}[style=python]
batch_size = 32  # Instead of 64
    \end{lstlisting}

    \item \textbf{Use gradient accumulation}:
    \begin{lstlisting}[style=python]
# Effective batch = 64, actual batch = 16
accumulation_steps = 4
batch_size = 16

for i, batch in enumerate(train_loader):
    loss = model(**batch)["loss"] / accumulation_steps
    loss.backward()

    if (i + 1) % accumulation_steps == 0:
        optimizer.step()
        optimizer.zero_grad()
    \end{lstlisting}

    \item \textbf{Reduce context length}:
    \begin{lstlisting}[style=python]
context_length = 256  # Instead of 512
    \end{lstlisting}

    \item \textbf{Enable mixed precision}:
    \begin{lstlisting}[style=python]
scaler = torch.cuda.amp.GradScaler()
with torch.cuda.amp.autocast():
    loss = model(**batch)["loss"]
    \end{lstlisting}

    \item \textbf{Clear cache}:
    \begin{lstlisting}[style=python]
torch.cuda.empty_cache()
    \end{lstlisting}
\end{enumerate}

\subsection{Memory Estimation}

\begin{equation}
\text{Memory} \approx 4 \times N + B \times L^2 \times H
\end{equation}

where:
\begin{itemize}
    \item $N$ = parameters (in float32 = 4 bytes each)
    \item $B$ = batch size
    \item $L$ = sequence length
    \item $H$ = number of heads (for attention matrices)
\end{itemize}

% -----------------------------------------------------------------------------
\section{Overfitting}
% -----------------------------------------------------------------------------

\textbf{Symptoms}: Training loss decreases but validation loss increases.

\subsection{Solutions}

\begin{enumerate}
    \item \textbf{Add dropout}:
    \begin{lstlisting}[style=python]
config = ModelConfig(dropout=0.2)  # Was 0.1
    \end{lstlisting}

    \item \textbf{Use more data}: TinyStories has 2.1M stories; ensure you're using all of it.

    \item \textbf{Reduce model size}:
    \begin{lstlisting}[style=python]
config = ModelConfig(
    n_layers=4,    # Was 6
    d_model=256,   # Was 384
)
    \end{lstlisting}

    \item \textbf{Add weight decay}:
    \begin{lstlisting}[style=python]
optimizer = AdamW(model.parameters(), weight_decay=0.1)
    \end{lstlisting}

    \item \textbf{Early stopping}:
    \begin{lstlisting}[style=python]
if val_loss > best_val_loss:
    patience_counter += 1
    if patience_counter >= patience:
        print("Early stopping!")
        break
    \end{lstlisting}
\end{enumerate}

% -----------------------------------------------------------------------------
\section{Repetitive Generation}
% -----------------------------------------------------------------------------

\textbf{Symptoms}: Model generates the same words or phrases repeatedly.

\subsection{Solutions}

\begin{enumerate}
    \item \textbf{Increase temperature}:
    \begin{lstlisting}[style=python]
temperature = 1.0  # Instead of 0.7
    \end{lstlisting}

    \item \textbf{Use repetition penalty}:
    \begin{lstlisting}[style=python]
def apply_repetition_penalty(logits, generated, penalty=1.2):
    for token in set(generated):
        logits[token] /= penalty
    return logits
    \end{lstlisting}

    \item \textbf{Use top-p sampling}:
    \begin{lstlisting}[style=python]
top_p = 0.9  # Nucleus sampling
    \end{lstlisting}

    \item \textbf{Train longer}: Model may not have learned enough diversity.
\end{enumerate}

% -----------------------------------------------------------------------------
\section{Slow Training}
% -----------------------------------------------------------------------------

\textbf{Symptoms}: Training takes longer than expected.

\subsection{Solutions}

\begin{enumerate}
    \item \textbf{Use torch.compile}:
    \begin{lstlisting}[style=python]
model = torch.compile(model)  # 1.5-2x speedup
    \end{lstlisting}

    \item \textbf{Enable mixed precision}:
    \begin{lstlisting}[style=python]
with torch.cuda.amp.autocast():
    loss = model(**batch)["loss"]
    \end{lstlisting}

    \item \textbf{Increase batch size}: Larger batches use GPU more efficiently.

    \item \textbf{Use multiple workers}:
    \begin{lstlisting}[style=python]
DataLoader(dataset, num_workers=4, pin_memory=True)
    \end{lstlisting}

    \item \textbf{Profile your code}:
    \begin{lstlisting}[style=python]
with torch.profiler.profile() as prof:
    model(**batch)
print(prof.key_averages().table())
    \end{lstlisting}
\end{enumerate}

% -----------------------------------------------------------------------------
\section{Debugging Checklist}
% -----------------------------------------------------------------------------

\begin{enumerate}
    \item[$\square$] Data loads correctly (decode and inspect samples)
    \item[$\square$] Labels are shifted by one position
    \item[$\square$] Attention mask matches padding
    \item[$\square$] Model parameters are on correct device
    \item[$\square$] Optimizer is updating parameters
    \item[$\square$] Learning rate schedule is correct
    \item[$\square$] Gradient clipping is enabled
    \item[$\square$] Loss is decreasing (not flat or NaN)
    \item[$\square$] Validation loss is tracked
    \item[$\square$] Checkpoints are being saved
\end{enumerate}

